%% Source: paper_evidence/mobile/iphone_15_pro_max_release_2026-11-20/
%%   console_phase_trace.txt (Part A) and
%%   extended_benchmark_1795167684106.json, memorySamples (Part B).
%% Rewritten 2026-11-20 (Release build); the earlier pre/post series had no surviving source.
\needspace{22\baselineskip}
\tablehead{tab:S7}{iOS resident memory with the contiguous output buffer on
the physical iPhone~15 Pro Max (one session, 20 November 2026, iOS~27.2.1,
committed release build). Values are mebibytes (MiB) of resident memory.}

\nopagebreak\noindent\textbf{Part A. Phase-boundary trace}

\begin{center}\footnotesize\setlength{\tabcolsep}{5pt}
\begin{tabular}{p{190pt} r}
\toprule
Phase boundary & Resident memory \\
\midrule
Start of benchmark                       & 34.1 \\
After cold-load loop                     & 34.7 \\
After models and benchmark images loaded & 105.0 \\
After classifier stage (110 calls)       & 141.6 \\
After detector stage, $1920\times1080$   & 205.7 \\
After detector stage, $1280\times960$    & 205.5 \\
After detector stage, $640\times480$     & 204.7 \\
After end-to-end loop (all five paths)   & 327.0 \\
\bottomrule
\end{tabular}\end{center}

\medskip
\noindent\textbf{Part B. Benchmark memory samples}

\begin{center}\footnotesize\setlength{\tabcolsep}{5pt}
\begin{tabular}{p{190pt} r}
\toprule
Sample & Resident memory \\
\midrule
Before the measured loops (after cold-load timing) & 34.9 \\
After all measured loops, including end-to-end     & 300.7 \\
\bottomrule
\end{tabular}\end{center}

{\footnotesize
\textit{Notes.} (1)~The application printed Part~A at each phase boundary;
the value is the process resident size divided by $1024^2$. Part~B's samples are written to the benchmark export at
different instants from the console lines, so the two parts need not agree.
(2)~Neither part samples inside a timed call, so a short-lived peak during a
single inference is not captured.
(3)~The earlier build passed each of the 226{,}800 detector output values to
shared code as a separate object. According to the session record of
14 September 2026, resident memory grew with every detector call until iOS
terminated the application at about 1{,}600~MiB during the first detector
stage. That process left no export, so no pre-correction trace is reported.
(4)~The correction replaced the boxed values with a contiguous numerical
buffer. Model weights, thresholds, class indexing, coordinate transformations
and non-maximum suppression were unchanged. An output-equivalence check in
the iOS Simulator reproduced all 226{,}800 values bit for bit.
(5)~No phase-boundary trace exists for Android. Android already used a
contiguous numerical representation. The Xiaomi memory difference is
excluded because both samples were captured at the same point.
\par}
\tabletail
